% Sección 9: Discusión, recomendaciones accionables y limitaciones
% evidencia:E16,E17,E18,E19,E20,E23,E27

\section{Discusión, recomendaciones accionables y limitaciones}

El valor práctico de MediFlow no radica en ofrecer cifras únicas de exactitud ni en automatizar decisiones clínicas, sino en articular datos administrativos heterogéneos bajo criterios cuantitativos auditables, dejando explícitos sus supuestos y márgenes de error. Una interpretación responsable exige contrastar los hallazgos numéricos frente a hipótesis alternativas y frente al estado real de implementación de cada componente.

\subsection{Interpretación operativa e hipótesis alternativas}

Los resultados predictivos confirman una frontera clara entre servicios de alto volumen y servicios de conteo bajo o intermitente. En Urgencias (8,4\,\% WAPE en \texttt{lgbm\_final}; 8,91\,\% en bundle \texttt{v2.0}) y Rehabilitación Ambulatoria entre lunes y viernes (85,6\,\%--90,9\,\% de exactitud), la combinación de rezagos semanales y calendario colombiano permite anticipar oscilaciones semanales y cierres festivos con error acotado. En contraste, en Tratamiento Urgencias ($\approx 15{,}2$ llegadas/día, WAPE 43,5\,\% en \texttt{lgbm\_final} y 42,18\,\% en \texttt{v2.0}) y en los fines de semana de Rehabilitación (exactitud de 51,3\,\% los sábados y 15,6\,\% los domingos), la métrica porcentual refleja tanto la sensibilidad de denominador pequeño como desplazamientos semanales de tasa por debajo del techo Poisson idealizado (77,0\,\% semanal observado frente a 92,3\,\% teórico en el experimento puntual) \parencite{mediflow2026demanda}. En consecuencia, la planificación en Tratamiento Urgencias debe apoyarse en el error absoluto medio (MAE $\approx 6{,}4$--$6{,}6$ pacientes/día) y en agregados semanales, evitando prometer alertas diarias precisas en porcentaje.

En el componente de congestión hospitalaria, la verificación archivada de Ley de Little ($277{,}52$ frente a $277{,}47$ pacientes medios en población alineada) confirma la consistencia matemática del censo reconstruido tras el recorte $p_{99{,}9}$ \parencite{mediflow2026informe}. Al mismo tiempo, dos resultados previenen interpretaciones simplistas. Primero, la correlación ecológica negativa ($r = -0{,}481$) entre el censo mensual de Hospitalización y la espera $p_{90}$ en la transición Tratamiento Urgencias $\rightarrow$ Hospitalización ($25{,}13$\,h en meses normales frente a $23{,}70$\,h en meses de censo alto) muestra que las demoras registradas entre marcas OMOP no responden linealmente a la saturación mensual agregada de piso; pueden obedecer a prácticas de cierre administrativo tardío, flujos por tipo de ingreso o solapamientos entre registros \parencite{mediflow2026optimizacion}. Segundo, la simulación de altas tempranas demuestra que desplazar el 35\,\% de egresos de la tarde a la mañana reduce las horas en estado crítico de 9,32\,\% a 6,26\,\% ($-32{,}8\,\%$ relativo; $-3{,}06$\,pp), pero deja intacto el máximo histórico de 359 pacientes. La descongestión modelada actúa sobre la duración acumulada de la presión, no sobre los picos extremos.

\subsection{Estado consolidado de capacidades implementadas y pendientes}

Para impedir que se confunda código versionado o ensayos sintéticos con capacidades desplegadas sobre datos hospitalarios reales, la Tabla~\ref{tab:estado_capacidades} sintetiza la situación exacta de cada componente al corte del commit \texttt{803c1a1}.

\begin{table}[htbp]
\caption{Estado verificado de capacidades entre código, artefacto servido y pendientes}
\label{tab:estado_capacidades}
{\singlespacing\small
\begin{tabularx}{\textwidth}{@{}L{0.22\textwidth} L{0.25\textwidth} L{0.24\textwidth} Y@{}}
\toprule
Capacidad o módulo & Implementación en código & Estado en bundle \texttt{v2.0} / PBIP & Trabajo pendiente \\
\midrule
Pronóstico puntual 1--7\,d (\texttt{lgbm\_final} Poisson) & Implementado en \texttt{demanda.py} con nivel local y recencia & No es el modelo servido; reporte archivado del \texttt{2026-11-18} & Evaluar incorporación de nivel local en familia cuantílica \\
Pronóstico cuantílico web ($p_{10}/p_{50}/p_{90}$ + SHAP) & Implementado en \texttt{explicacion.py} y \texttt{predict.py} & Activo en bundle \texttt{v2.0} (corte \texttt{2026-07-17}) & Recalibrar cobertura (64,3\,\%--78,6\,\%) e incluir Carnaval/Semana Santa \\
Largo plazo a 120 días & Implementado en \texttt{largo\_plazo.py} (banda base $p_{05}/p_{95}$ ensanchada) & Activo vía \texttt{/api/largo-plazo} (corrida propia en \texttt{v2.0}) & Reconciliar reporte del \texttt{2026-11-03} frente a corrida empaquetada \\
Erlang~C histórico, Little y altas tempranas & Implementado en \texttt{optimizacion.py} (interpolación en API) & Activo en \texttt{v2.0} con Little híbrido (\num{284,01} pac.) & Regenerar artefacto con Little alineado (\num{277,52}) y medir $1/\mu$ real \\
Acoplamiento por corte Etapa~1 (\texttt{perfiles\_horarios}) & Implementado en \texttt{empaquetar.py} (\texttt{v2.1}) y probado con fixtures & No habilitado en \texttt{v2.0}: devuelve \texttt{disponible: false} & Regenerar paquete \texttt{v2.1} autorizado con perfiles sin fuga \\
Asistente \texttt{/api/chat} & 6 herramientas deterministas + respaldo offline (\texttt{chat.py}) & Operativo contra salidas del bundle vía FastAPI & Validar funcionamiento end-to-end con clave autorizada en despliegue \\
Tablero Power BI (PBIP) & 7 pág., 12 tab., 18 rel., 43 medidas DAX (\texttt{FechaCorte = 2026-04-30}) & Solo definiciones locales corregidas; sin datos en PBIP & Revisar en Desktop, alinear corte abril vs. julio y republicar enlace \\
Simulador de recorrido (Etapa~2) y piloto en sombra (Etapa~3) & Diseño conceptual documentado en \texttt{informe\_final.md} & No implementado en software ni validado en clínica & Instrumentar 3 marcas HIS, capacidad real y aprobar piloto \\
\bottomrule
\end{tabularx}}
\apanota{Consolidación basada en \textcite{mediflow2026auditoria}, \textcite{mediflow2026bundle}, \textcite{mediflow2026pbip} y la inspección directa de \texttt{src/hun/models/} y \texttt{app/}.}
\end{table}

\subsection{Recomendaciones condicionadas y hoja de ruta evolutiva}

Las cinco recomendaciones formuladas en la Tabla~\ref{tab:recomendaciones} (\texttt{REC-01} a \texttt{REC-05}) sintetizan los escenarios cuantificados del sistema \parencite{mediflow2026optimizacion,mediflow2026informe}. Todas tienen carácter de propuestas condicionadas a validación institucional previa; no constituyen compromisos acordados con las jefaturas del HUN ni proyecciones de ahorro económico o dotación contractual.

\begin{table}[htbp]
\caption{Recomendaciones operacionales condicionadas (\texttt{REC-01} a \texttt{REC-05}) y criterios de validación}
\label{tab:recomendaciones}
{\singlespacing\small
\begin{tabularx}{\textwidth}{@{}l L{0.24\textwidth} L{0.30\textwidth} Y@{}}
\toprule
ID & Acción propuesta & \shortstack[l]{Resultado modelado\\o archivado} & Condición e indicador de piloto \\
\midrule
\texttt{REC-01} & Validar equipos equivalentes en Urgencias (06--12\,h) & A demanda $p_{90}$ (39,67/h): 15 equipos frente a 12 normales; espera media 6,09\,min & Medir duración clínica real frente al supuesto de 20\,min y plantilla instalada por turno \\
\texttt{REC-02} & Revisar capacidad de Rehabilitación (06--12\,h) & A demanda $p_{90}$ (133,33/h): 70 equipos frente a 42 normales; espera media 5,30\,min & Contrastar agenda programada, atención supuesta (30\,min) y puestos simultáneos \\
\texttt{REC-03} & Investigar demoras evitables en estancia de Hospitalización & Reducir 4\,h libera teóricamente 13,97 camas medias equivalentes (4,92\,\% del Little híbrido) & Auditar causas administrativas de egreso manteniendo intactos los criterios clínicos de alta \\
\texttt{REC-04} & Pilotar egreso seguro intradía en altas elegibles & Adelantar 35\,\% de [14,20) a [08,12): horas rojas 9,32\,\% $\rightarrow$ 6,26\,\% ($-32{,}8\,\%$ rel.); pico 359 & Evaluar en piloto con umbral fijo ($>309$), tasa de reingreso y seguridad del alta \\
\texttt{REC-05} & Auditar sellos temporales en Tratamiento Urg. $\rightarrow$ Hosp. & \num{47670} pasos; $p_{90}$ mensual ponderado 24,45\,h; correlación con censo $r = -0{,}481$ & Instrumentar decisión, asignación de cama y traslado efectivo antes de atribuir demora a cama \\
\bottomrule
\end{tabularx}}
\apanota{Propuestas extraídas de \textcite{mediflow2026optimizacion} y \textcite{mediflow2026informe}. Los porcentajes y camas medias equivalentes representan escenarios retrospectivos o paramétricos; los indicadores sugeridos para una eventual fase piloto (demora efectiva de traslado, estancia observada, horas en alerta roja a umbral fijo y métricas de seguridad clínica) no implican metas institucionales aprobadas ni calendario confirmado.}
\end{table}

La hoja de ruta para trasladar estas propuestas a operación asistencial se estructura en tres etapas secuenciales \parencite{mediflow2026informe}: (1) completar el paquete \texttt{v2.1} sobre agregados autorizados incorporando perfiles horarios por corte, reconciliando el valor de Little con población común y recalibrando cuantiles por servicio; (2) desarrollar el simulador de recorrido integrando camas habilitadas, dotación efectiva, tiempos de atención medidos y tres sellos del HIS (\texttt{hora\_decision\_hospitalizar}, \texttt{hora\_asignacion\_cama} y \texttt{hora\_traslado\_efectivo}); y (3) ejecutar un piloto prospectivo en modo sombra paralelo a la operación habitual antes de validar cualquier intervención asistencial supervisada.

\subsection{Limitaciones consolidadas}

Seis limitaciones estructurales acotan el uso actual de MediFlow: (a) ausencia en OMOP de camas físicas habilitadas, dotación de personal y tiempos clínicos de atención, obligando a trabajar con equipos equivalentes y umbrales históricos relativos; (b) carácter estadístico del recorte $p_{99{,}9}$ sobre estancias extremas y exclusión de variables clínicas de diagnóstico, procedimiento o laboratorio en la matriz predictiva; (c) ambigüedad administrativa de \texttt{espera\_h} en transiciones con posibles solapamientos negativos; (d) cobertura desigual en los intervalos cuantílicos del bundle \texttt{v2.0} (inferiores a 80\,\% nominal en los tres servicios sobre 182 días) y ausencia de perfiles por corte en el binario actual; (e) divergencia temporal entre Power BI (corte abril 2026 sin verificación de republicación) y demanda (corte julio 2026), así como entre los dos reportes comprobados de largo plazo; y (f) inexistencia de validación prospectiva en vivo o autorización clínica para decisiones asistenciales directas.
